\documentclass[twocolumn,10pt,a4paper]{article}
\usepackage[T1]{fontenc}
\usepackage[utf8]{inputenc}
\usepackage{mathptmx}
\usepackage{amsmath}
\usepackage[a4paper,top=0.95in,bottom=1.0in,left=0.78in,right=0.78in,columnsep=0.30in]{geometry}
\usepackage{graphicx}
\usepackage{booktabs}
\usepackage{array}
\usepackage{caption}
\usepackage{authblk}
\usepackage{titlesec}
\usepackage{enumitem}
\usepackage{fancyhdr}
\usepackage{xcolor}
\usepackage[colorlinks=true,linkcolor=linknavy,citecolor=linknavy,urlcolor=linknavy]{hyperref}

\definecolor{linknavy}{HTML}{1B4F72}
\definecolor{rule}{HTML}{BBBBBB}

% --- section styling, REVTeX-like small caps roman numerals ---
\renewcommand{\thesection}{\Roman{section}}
\renewcommand{\thesubsection}{\thesection\,\Alph{subsection}}
\titleformat{\section}{\centering\normalsize\scshape}{\thesection.}{0.6em}{}
\titleformat{\subsection}{\normalsize\itshape}{\Alph{subsection}.}{0.5em}{}
\titlespacing*{\section}{0pt}{11pt}{5pt}
\titlespacing*{\subsection}{0pt}{8pt}{3pt}

\captionsetup{font=small,labelfont=bf,labelsep=period,justification=justified,
              singlelinecheck=false,skip=5pt}
\captionsetup[table]{position=top,skip=4pt}

\setlist[itemize]{leftmargin=1.1em,itemsep=1.5pt,topsep=3pt,parsep=0pt}
\setlist[enumerate]{leftmargin=1.5em,itemsep=1.5pt,topsep=3pt,parsep=0pt}

\renewcommand{\arraystretch}{1.12}
\setlength{\parindent}{1.2em}
\setlength{\skip\footins}{10pt}
\renewcommand{\footnoterule}{\noindent\rule{0.4\columnwidth}{0.4pt}\vspace{3pt}}

\pagestyle{fancy}\fancyhf{}
\renewcommand{\headrulewidth}{0pt}
\fancyfoot[C]{\small\thepage}
\fancypagestyle{titlestyle}{\fancyhf{}\renewcommand{\headrulewidth}{0pt}%
  \fancyfoot[C]{\small\thepage}}

\newcommand{\RBI}{\textrm{RBI}}
\newcommand{\IARS}{\textrm{IARS}}
\newcommand{\SCS}{\textrm{SCS}}
\newcommand{\CRS}{\textrm{CRS}}

% --- small ORCID mark, drawn rather than loaded as an image ---
\usepackage{tikz}
\newcommand{\orcidmark}{\raisebox{0.15ex}{\begin{tikzpicture}[scale=0.055]
  \fill[color=green!55!black] (0,0) circle (1);
  \fill[white] (-0.33,-0.52) rectangle (-0.15,0.30);
  \fill[white] (-0.24,0.52) circle (0.12);
  \draw[white,line width=0.18] (0.05,-0.45) -- (0.05,0.33);
  \draw[white,line width=0.18] (0.05,0.33) .. controls (0.75,0.33) and (0.75,-0.45) .. (0.05,-0.45);
\end{tikzpicture}}}


\begin{document}

\twocolumn[
\begin{@twocolumnfalse}
\vspace{-6mm}
\begin{center}
{\large\bfseries Procedure-level biorisk scoring for aerosol-generating bench steps\\[2pt]
in clinical diagnostic laboratories: an open instrument,\\[2pt]
reference scores, and an evidence map\par}
\vspace{5mm}
{\normalsize Oyekunle Adedamola Oyeniyi\,\orcidmark\par}
\vspace{1.5mm}
{\small\itshape Independent Researcher, Pullman, Washington 99163, USA\par}
\vspace{1mm}
{\footnotesize ORCID: \href{https://orcid.org/0000-0002-1978-6133}{0000-0002-1978-6133}
\quad$\cdot$\quad \href{mailto:adedamolaoyekunleniyi@gmail.com}{adedamolaoyekunleniyi@gmail.com}\par}
\vspace{2mm}
{\footnotesize Instrument version 1.1 \quad$\cdot$\quad 29 October 2026\par}
\end{center}
\vspace{4mm}


\noindent\hspace{0.055\textwidth}\begin{minipage}{0.89\textwidth}
\small
\noindent Clinical laboratories are required by national biosafety guidance to assess
the biorisk of their own procedures, but no openly available instrument operates at the
level of an individual bench step, and the evidence underlying routine bench biosafety
practice has never been systematically characterised. I developed a procedure-level
scoring instrument of twelve anchored items across three dimensions --- inherent aerosol
release, scenario consequence, and control reliability --- with an evidence grade reported
beside every score and never combined into it, and mapped each element to the clause
structure of ISO 35001:2019. I applied the instrument to seven routine bench steps under
three agent profiles, extracting findings from openly accessible measurement studies,
incident investigations, and guidance. Across seventeen step-and-profile combinations,
residual biorisk under baseline practice ranged from 3.33 to 19.60 on a 1--25 scale, and
the recommended control sets reduced it by a mean of 4.48 points. Direct quantitative
aerosol measurement existed for three of seven steps. The ordering that resulted is
inverted with respect to the evidence: the two steps carrying the highest residual
biorisk, open-tube culture work (19.60, ``Very high'') and MALDI-TOF target preparation
(17.73, ``High''), rest on incident evidence with no aerosol measurement of either step in
a clinical laboratory, while the pre-analytical steps most often cited in requests for
containment --- tube opening and centrifugation --- have been measured directly and
yielded no recoverable organisms. The highest-throughput step in the laboratory, operation
of automated chemistry and haematology analysers, carries no located evidence of any kind.
I argue that this inversion is a property of how the evidence base was assembled rather
than of the hazards themselves, that the negative pre-analytical results are evidence for
the efficacy of containment rather than for the adequacy of the open bench, and that four
specific measurement studies would resolve most of the uncertainty documented here. The
instrument, the scores, the extracted evidence, and the analysis code are released openly.

\vspace{2.5mm}
\noindent\textbf{Keywords:} biorisk management; clinical laboratory; infectious aerosol;
laboratory-acquired infection; ISO 35001; biosafety cabinet; evidence grading; sentinel
laboratory
\end{minipage}

\vspace{7mm}
\end{@twocolumnfalse}
]
\thispagestyle{titlestyle}
